\input{preamble}

\begin{document}

\title{Economy}
\maketitle

\phantomsection
\label{section-phantom}

\tableofcontents

\section{What energy actually means}
\label{section-energy-means}

\noindent
This is an AI-generated section.

\medskip\noindent
This section follows the historical
framework developed by Vaclav Smil in
\emph{Energy and Civilization: A History}
\cite{smil2017energy}.

\medskip\noindent
{\bf Energy and useful work.}

Energy is not merely one commodity among
others. Every society depends on converting
energy into useful work: movement, heat,
light, construction, transportation, food
production, and, more recently,
computation. What changes historically is
not the existence of energy, but the ability
of human societies to capture, convert,
store, transport, and control energy flows.

\medskip\noindent
{\bf Foraging societies.}

For early humans, almost all available
energy ultimately came from photosynthesis
[but I do think that all the energy on Earth
ultimately comes from the Sun,
because, isn't oil coming from
very old fossilized organic matter?].
Plants stored solar energy, humans and
animals consumed plants or other animals,
and muscular effort converted food energy
into mechanical work.

Fire made it possible to obtain heat outside
the human body, while tools increased the
effectiveness of muscular labour. Still,
the amount of useful work available to a
society remained closely constrained by
human metabolism and by the biological
productivity of its environment.

\medskip\noindent
{\bf Agriculture and the control of solar
energy.}

Agriculture reorganized ecosystems in order
to concentrate solar energy into crops and
domesticated animals useful to humans. It
therefore represented much more than a new
method of obtaining food.

Agricultural surpluses supported larger
populations, permanent settlements,
specialized occupations, armies, and
states. Nevertheless, traditional
agricultural civilizations remained
dependent on current or recent solar flows:
food, wood, animal feed, wind, and flowing
water.

Their energy supply was consequently
limited by land, climate, seasons, and the
rate at which biological resources could
be renewed.

\medskip\noindent
{\bf Animals, water, and wind.}

Draft animals supplied mechanical power
beyond the human body. Waterwheels and
windmills converted natural flows into
useful work without requiring continuous
human or animal effort.

These sources considerably increased the
power available for farming, milling,
transport, and manufacturing. They were,
however, geographically and temporally
restricted. Rivers existed only in certain
places, winds were variable, and animals
required large quantities of land and feed.

\medskip\noindent
{\bf Fossil fuels and industrialization.}

Coal, and later oil and natural gas, allowed
societies to use energy accumulated over
geological timescales. Humanity was no
longer restricted to the energy captured
each year by contemporary biological
processes.

Fossil fuels were particularly
transformative because they were highly
concentrated, storable, transportable, and
capable of powering machines on demand.
Steam engines converted coal into
mechanical work, while internal-combustion
engines made petroleum central to modern
transportation.

The resulting increase in available power
transformed industry, agriculture,
transport, warfare, cities, and ordinary
domestic life.

\medskip\noindent
{\bf Electricity as an energy carrier.}

Electricity is not a primary source of
energy. It must be produced from some other
source, such as coal, natural gas, moving
water, nuclear reactions, wind, or
sunlight.

Its importance lies in its flexibility.
Electricity can be transported through
networks and converted into light, heat,
motion, refrigeration, communication, or
computation. Modern societies therefore
depend not only on producing energy, but
also on maintaining large and precisely
coordinated electrical systems.

\medskip\noindent
{\bf The historical lesson.}

Smil's central perspective is that the
history of civilization can partly be read
as the history of increasingly powerful
energy conversions:
\[
  \text{capture}
  \longrightarrow
  \text{conversion}
  \longrightarrow
  \text{control}
  \longrightarrow
  \text{useful work}.
\]

Civilizations are not shaped merely by the
total amount of energy they possess. They
are shaped by the concentration, quality,
reliability, location, and convertibility
of their energy sources.

This provides the material background for
the economic questions studied later.
Markets, governments, and institutions
organize who controls energy resources,
how they are used, who receives their
benefits, and who bears the environmental
and social costs.

\bibliography{my}
\bibliographystyle{amsalpha}

\end{document}
